\section{Результаты}\label{sec:results}

\subsection{Покрытие при дрейфе позы}\label{sec:coverage}

\begin{table}[H]
\centering\scriptsize
\caption{Покрытие гарантии и медианный запас $\hat r$ по 200 разбиениям, $1 - \alpha = 0{,}9$. Медианная ATE уровней: L2 --- 2~см, L3 --- 7~см, L4 --- 45~см. «отк.» --- доля разбиений с отказом, если она не меньше половины}
\label{tab:coverage}
\begin{tabular}{llcccccccc}
\toprule
Класс & Калибровка & \multicolumn{2}{c}{L0} & \multicolumn{2}{c}{L2} & \multicolumn{2}{c}{L3} & \multicolumn{2}{c}{L4} \\
 & & покр. & $\hat r$, м & покр. & $\hat r$, м & покр. & $\hat r$, м & покр. & $\hat r$, м \\
\midrule
велосипед & без калибровки & 0{,}00 & 0 & 0{,}00 & 0 & 0{,}01 & 0 & 0{,}01 & 0 \\
 & метки клеток & 1{,}00 & 0 & 0{,}73 & 0 & 0{,}34 & 0 & 0{,}17 & 0 \\
 & только метки & 0{,}93 & 0{,}75 & 0{,}94 & 0{,}75 & 0{,}93 & 0{,}75 & 0{,}76 & 0{,}75 \\
 & только поза & 0{,}00 & 0{,}00 & 0{,}03 & 0{,}07 & 0{,}22 & 0{,}18 & 0{,}82 & 0{,}90 \\
 & раздельно & 0{,}93 & 0{,}75 & 0{,}98 & 0{,}82 & 0{,}99 & 0{,}90 & 0{,}97 & 1{,}61 \\
 & \textbf{совместно} & 0{,}93 & 0{,}75 & 0{,}94 & 0{,}75 & 0{,}94 & 0{,}75 & 0{,}93 & 1{,}18 \\
\midrule
растение & без калибровки & 0{,}00 & 0 & 0{,}00 & 0 & 0{,}00 & 0 & 0{,}00 & 0 \\
 & метки клеток & 1{,}00 & 0 & 0{,}41 & 0 & 0{,}10 & 0 & 0{,}03 & 0 \\
 & только метки & 0{,}93 & отк.\,0{,}63 & 0{,}96 & отк.\,0{,}64 & 0{,}94 & отк.\,0{,}64 & 0{,}86 & отк.\,0{,}64 \\
 & только поза & 0{,}00 & 0{,}00 & 0{,}00 & 0{,}11 & 0{,}13 & 0{,}25 & 0{,}86 & 1{,}39 \\
 & раздельно & 0{,}93 & отк.\,0{,}63 & 0{,}98 & отк.\,0{,}64 & 0{,}96 & отк.\,0{,}64 & 0{,}95 & отк.\,0{,}64 \\
 & \textbf{совместно} & 0{,}93 & отк.\,0{,}63 & 0{,}93 & 0{,}80 & 0{,}94 & отк.\,0{,}55 & 0{,}93 & отк.\,0{,}58 \\
\midrule
тумба под ТВ & без калибровки & 0{,}05 & 0 & 0{,}01 & 0 & 0{,}00 & 0 & 0{,}00 & 0 \\
 & метки клеток & 1{,}00 & 0 & 0{,}77 & 0 & 0{,}61 & 0 & 0{,}30 & 0 \\
 & только метки & 0{,}98 & отк.\,0{,}95 & 0{,}98 & отк.\,0{,}95 & 0{,}98 & отк.\,0{,}95 & 0{,}99 & отк.\,0{,}95 \\
 & только поза & 0{,}05 & 0{,}00 & 0{,}06 & 0{,}07 & 0{,}13 & 0{,}16 & 0{,}47 & 0{,}56 \\
 & раздельно & 0{,}98 & отк.\,0{,}95 & 0{,}98 & отк.\,0{,}95 & 0{,}98 & отк.\,0{,}95 & 1{,}00 & отк.\,0{,}95 \\
 & \textbf{совместно} & 0{,}98 & отк.\,0{,}95 & 0{,}98 & отк.\,0{,}94 & 0{,}97 & отк.\,0{,}93 & 0{,}94 & отк.\,0{,}72 \\
\bottomrule
\end{tabular}

\end{table}

\begin{figure}[H]
\centering
\includegraphics[width=\linewidth]{figures/coverage_headline.png}
\caption{Покрытие (слева) и запас (справа), среднее по растению и велосипеду}
\label{fig:coverage}
\end{figure}

Табл.~\ref{tab:coverage} и рис.~\ref{fig:coverage} проверяют гипотезы H1--H3. Карта без калибровки почти никогда не накрывает истинный объект целиком: детектор обрезает края.

Калибровка меток клеток даёт покрытие 1,00 без дрейфа, но лишь потому, что её порог $\lambda^\ast = 0$ объявляет возможным растением любую занятую клетку. Запаса в метрах у неё нет, поэтому сдвиг на 2~см уже снижает покрытие велосипеда до 0,73.

Запас «только метки» держится, пока дрейф мал по сравнению с ошибкой меток (до 7~см), и падает до 0,76 при 45~см (H1). Запас «только поза» не видит ошибок детектора: при дрейфе до 7~см его покрытие не выше 0,22 (H2).

Совместный запас держит 0,93--0,94 на всех уровнях. При 45~см он равен 1,18~м для велосипеда против 1,61~м у суммы отдельных запасов и 1,53 против 2,11~м для растения, то есть на 27--28\,\% меньше (H3).

Растение отказывается примерно в 60\,\% разбиений: в одной сцене детектор пропускает его целиком, а при 13 сценах запас равен максимуму. Тумба под ТВ отказывается в 71--95\,\% разбиений: в двух сценах из 21 детектор уверенно принимает её за скамейку или ковёр (промах 5,6 и 6,5~м). Для этого детектора гарантировать тумбу на уровне 90\,\% нельзя, и калибровка сообщает об этом.

\subsection{Миссии}\label{sec:missions}

\begin{figure}[H]
\centering
\includegraphics[width=0.7\linewidth]{figures/example_v3_sc0_staging_20_paper.png}
\caption{Трудная задача при дрейфе L2. Красное --- истинные растение и велосипед, серое --- препятствия, линия --- граница запретной зоны. Без калибровки зона не накрывает часть растения, и путь проходит ближе 0,5~м (чёрные точки); совместный запас накрывает объект}
\label{fig:example}
\end{figure}

\begin{table}[H]
\centering\small
\caption{Трудные задачи (растение и велосипед, 233 задачи): доля задач с найденным путём, доля опасных среди найденных путей, доля успешных миссий}
\label{tab:missions}
\begin{tabular}{lccccccccc}
\toprule
Ветка & \multicolumn{3}{c}{L0} & \multicolumn{3}{c}{L2} & \multicolumn{3}{c}{L4} \\
 & путь & наруш. & успех & путь & наруш. & успех & путь & наруш. & успех \\
\midrule
без калибровки & 0{,}99 & 0{,}27 & 0{,}73 & 0{,}97 & 0{,}32 & 0{,}65 & 0{,}26 & 0{,}34 & 0{,}17 \\
покадровая CP меток & 0{,}27 & 0{,}00 & 0{,}27 & 0{,}28 & 0{,}00 & 0{,}28 & 0{,}04 & 0{,}23 & 0{,}03 \\
только метки & 0{,}32 & 0{,}00 & 0{,}32 & 0{,}31 & 0{,}00 & 0{,}31 & 0{,}10 & 0{,}00 & 0{,}10 \\
только поза & 0{,}99 & 0{,}27 & 0{,}73 & 0{,}96 & 0{,}20 & 0{,}77 & 0{,}09 & 0{,}00 & 0{,}09 \\
раздельно (сумма) & 0{,}32 & 0{,}00 & 0{,}32 & 0{,}31 & 0{,}00 & 0{,}31 & 0{,}00 & 0{,}00 & 0{,}00 \\
\textbf{совместно} & 0{,}32 & 0{,}00 & 0{,}32 & 0{,}31 & 0{,}00 & 0{,}31 & 0{,}03 & 0{,}00 & 0{,}03 \\
оракул & 1{,}00 & 0{,}00 & 1{,}00 & 0{,}98 & 0{,}00 & 0{,}98 & 0{,}98 & 0{,}00 & 0{,}98 \\
\bottomrule
\end{tabular}

\end{table}

На трудных задачах планировщик без калибровки находит путь почти всегда, но 27--37\,\% его путей опасны (L0--L3). У всех вариантов с откалиброванным запасом опасных путей нет. Цена --- путь находится лишь в 30--32\,\% трудных задач, тогда как на истинной карте --- в 98--100\,\%. На случайных задачах совместный запас даёт 78\,\% успешных миссий без дрейфа и 68\,\% при дрейфе L3, без опасных путей. Без калибровки успешны 95 и 82\,\%, но 4--6\,\% миссий опасны.

Если запретными считаются все три класса, отказ тумбы переводит планировщик на запасной вариант, и путь находится лишь в 23\,\% трудных задач без дрейфа и в 5--8\,\% при дрейфе L2--L3.

\subsection{Сравнение с калибровкой меток в её условиях}\label{sec:comparison}

Успех миссий у~\cite{sundarsingh} (98,36\,\%) получен в другой постановке: открытые площадки, известная поза, пять классов, случайные задачи и режим исследования, когда пути нет. Поэтому мы сравниваем методы на наших сценах в их условиях. Поза известна. Словарь закрыт пятью классами: из 98 текстовых эмбеддингов детектора оставлены пять, что совпадает с перенастройкой словаря. Цель $1 - \alpha = 0{,}95$, задачи случайные. Режим исследования не воспроизводился.

\begin{table}[H]
\centering\scriptsize\setlength{\tabcolsep}{3.5pt}
\caption{Случайные задачи: покрытие (минимум по растению и велосипеду), доля успешных миссий и доля опасных миссий}
\label{tab:comparison}
\begin{tabular}{ll|ccc|ccc|cc|c}
\toprule
Условия & Дрейф & \multicolumn{3}{c|}{метки клеток} & \multicolumn{3}{c|}{\textbf{совместно}} & \multicolumn{2}{c|}{без калибровки} & оракул \\
 & & покр. & успех & опасно & покр. & успех & опасно & успех & опасно & успех \\
\midrule
закрытый словарь, 95\,\% & известна & 1{,}00 & 0{,}43 & 0{,}000 & 0{,}94 & 0{,}42 & 0{,}000 & 0{,}96 & 0{,}040 & 1{,}00 \\
 & L2 & 0{,}42 & 0{,}40 & 0{,}000 & 0{,}94 & 0{,}61 & 0{,}000 & 0{,}92 & 0{,}042 & 0{,}97 \\
 & L3 & 0{,}10 & 0{,}32 & 0{,}002 & 0{,}94 & 0{,}50 & 0{,}000 & 0{,}82 & 0{,}049 & 0{,}97 \\
\midrule
открытый словарь, 90\,\% & известна & 1{,}00 & 0{,}43 & 0{,}000 & 0{,}93 & 0{,}78 & 0{,}000 & 0{,}95 & 0{,}043 & 1{,}00 \\
 & L2 & 0{,}41 & 0{,}40 & 0{,}000 & 0{,}93 & 0{,}74 & 0{,}000 & 0{,}91 & 0{,}050 & 0{,}97 \\
 & L3 & 0{,}10 & 0{,}32 & 0{,}002 & 0{,}94 & 0{,}68 & 0{,}000 & 0{,}82 & 0{,}056 & 0{,}97 \\
 & L4 & 0{,}03 & 0{,}09 & 0{,}000 & 0{,}93 & 0{,}17 & 0{,}000 & 0{,}29 & 0{,}021 & 0{,}98 \\
\bottomrule
\end{tabular}

\end{table}

\begin{figure}[H]
\centering
\includegraphics[width=\linewidth]{figures/comparison.png}
\caption{Калибровка меток клеток~\cite{sundarsingh} и совместный запас на случайных задачах. Сверху --- закрытый словарь и $1 - \alpha = 0{,}95$, снизу --- открытый словарь и $1 - \alpha = 0{,}9$}
\label{fig:comparison}
\end{figure}

В их условиях методы равны: 0,43 и 0,42 успешных миссий, опасных нет (табл.~\ref{tab:comparison}). При дрейфе 2 и 7~см покрытие калибровки меток падает до 0,42 и 0,10, совместный запас держит 0,94 и даёт больше успешных миссий: 0,61 и 0,50 против 0,40 и 0,32. С открытым словарём совместный запас даёт в 1,8 раза больше успешных миссий уже при известной позе.

Калибровка меток у нас вырождается: порог $\lambda^\ast = 0$ в обоих словарях. Мы проверили, не виновата ли мелкая сетка, и повторили обе калибровки на сетках 5, 25, 50 и 100~см. Порог равен нулю на всех; запретная зона калибровки меток занимает 64\,\% свободного пространства при 5~см и всё пространство при 1~м. Причина в полных промахах детектора: если в одной калибровочной сцене до части объекта не дошло ни одного наблюдения его класса, худшая клетка даёт счёт 1 и порог 0. Множество меток может сказать только «эта клетка может быть растением». Расстояние говорит «растение где-то рядом», и это остаётся полезным, когда детектор теряет части объектов.

\subsection{Цена гарантии}\label{sec:price}

\begin{table}[H]
\centering\small
\caption{Уровень риска $\alpha$ при дрейфе L3: покрытие и запас совместного варианта, доля отказов для растения, доля трудных задач с найденным путём и доля опасных миссий. При $\alpha = 0{,}05$ разбиение 19/2}
\label{tab:alpha}
\begin{tabular}{c|cc|cc|c|cc|cc}
\toprule
$\alpha$ & \multicolumn{2}{c|}{покрытие} & \multicolumn{2}{c|}{$\hat r$, м} & отказ & \multicolumn{2}{c|}{трудные: путь} & \multicolumn{2}{c}{случайные: успех} \\
 & раст. & велос. & раст. & велос. & раст. & совм. & оракул & совм. & оракул \\
\midrule
0{,}05 & 0{,}95 & 0{,}95 & 0{,}85 & 0{,}76 & 0{,}70 & 0{,}10 & 0{,}99 & 0{,}34 & 0{,}97 \\
0{,}10 & 0{,}93 & 0{,}93 & 0{,}80 & 0{,}75 & 0{,}52 & 0{,}30 & 0{,}99 & 0{,}68 & 0{,}97 \\
0{,}20 & 0{,}86 & 0{,}86 & 0{,}74 & 0{,}70 & 0{,}00 & 0{,}31 & 0{,}99 & --- & --- \\
0{,}30 & 0{,}72 & 0{,}72 & 0{,}60 & 0{,}50 & 0{,}00 & 0{,}31 & 0{,}99 & 0{,}71 & 0{,}97 \\
\bottomrule
\end{tabular}

\end{table}

Повышение риска почти не добавляет решаемых задач (табл.~\ref{tab:alpha}). С $\alpha = 0{,}1$ до 0,3 запас уменьшается на 0,2--0,25~м, а доля трудных задач с путём растёт с 30 до 31\,\%; покрытие при этом следует за $1 - \alpha$.

Лучшие маски тоже не уменьшают запас. MobileSAM~\cite{mobilesam}, получающий на вход рамки детектора, снижает медианный промах велосипеда с 0,43 до 0,05~м. Однако наихудшая сцена сохраняет промах 0,75~м, и при 13 калибровочных сценах запас равен именно ей. С масками SAM карта без калибровки безопасна при известной позе (опасных путей нет), но не при дрейфе (9,5\,\% опасных путей при L3).

Запас определяет хвост ошибок восприятия. Уменьшить его можно, убрав наихудшие ошибки детектора или взяв больше калибровочных сцен.

\subsection{Когда сумма запасов недопокрывает}\label{sec:composition}

На синтетическом дрейфе сумма отдельных запасов только переплачивает: дрейф независим от ошибок детектора. Чтобы проверить утверждение~\ref{prop:sep}, мы оставили реальные ошибки меток и заменили дрейф сценовой смесью. В $k$ сценах из 21 происходит сбой локализации (реализации уровня L5), в остальных дрейф пренебрежимо мал (L1). Сцены со сбоем выбираются там, где метки хуже всего, где они лучше всего, или случайно. Уровень $\alpha = 0{,}3$ выбран потому, что при $\alpha = 0{,}1$ и 13 сценах любой сбой попадает в максимум.

\begin{figure}[H]
\centering
\includegraphics[width=\linewidth]{figures/composition_stress_a03_L5.png}
\caption{Покрытие суммы отдельных запасов (розовое) и совместного запаса (синее) в зависимости от доли сцен со сбоем локализации; $1 - \alpha = 0{,}7$, нижняя линия --- граница $1 - 2\alpha$}
\label{fig:composition}
\end{figure}

Если сбои совпадают с плохими метками, сумма даёт 0,76--0,84. Если сбои приходятся на сцены с хорошими метками, она опускается до 0,66 при цели 0,70 (растение, 19--24\,\% сцен со сбоем). Совместный запас даёт 0,71--0,78 при любом расположении сбоев (рис.~\ref{fig:composition}).

\subsection{Реальная одометрия и другие сцены}\label{sec:variants}

\begin{table}[H]
\centering\scriptsize\setlength{\tabcolsep}{3.5pt}
\caption{Проверка упрощений: покрытие и медианный запас (м), $\alpha = 0{,}1$; $^\dagger$ --- вариант отказывается в половине разбиений или чаще, приведён запас в остальных}
\label{tab:variants}
\begin{tabular}{l|cccccc|cccccc}
\toprule
Вариант & \multicolumn{6}{c|}{велосипед} & \multicolumn{6}{c}{растение} \\
 & \multicolumn{2}{c}{метки} & \multicolumn{2}{c}{раздельно} & \multicolumn{2}{c}{\textbf{совместно}} & \multicolumn{2}{c}{метки} & \multicolumn{2}{c}{раздельно} & \multicolumn{2}{c}{\textbf{совместно}} \\
 & покр. & $\hat q$ & покр. & $\hat q$ & покр. & $\hat q$ & покр. & $\hat q$ & покр. & $\hat q$ & покр. & $\hat q$ \\
\midrule
рамка + глубина, L0 & 0{,}93 & 0{,}75 & 0{,}93 & 0{,}75 & 0{,}93 & 0{,}75 & 0{,}93 & 0{,}85$^\dagger$ & 0{,}93 & 0{,}85$^\dagger$ & 0{,}93 & 0{,}85$^\dagger$ \\
рамка + глубина, L3 & 0{,}93 & 0{,}75 & 0{,}99 & 0{,}90 & 0{,}94 & 0{,}75 & 0{,}94 & 0{,}85$^\dagger$ & 0{,}96 & 1{,}09$^\dagger$ & 0{,}94 & 0{,}81$^\dagger$ \\
рамка + глубина, L4 & 0{,}76 & 0{,}75 & 0{,}97 & 1{,}61 & 0{,}93 & 1{,}18 & 0{,}86 & 0{,}85$^\dagger$ & 0{,}95 & 2{,}11$^\dagger$ & 0{,}93 & 1{,}53$^\dagger$ \\
MobileSAM, L0 & 0{,}93 & 0{,}75 & 0{,}93 & 0{,}75 & 0{,}93 & 0{,}75 & 0{,}93 & 0{,}61$^\dagger$ & 0{,}93 & 0{,}61$^\dagger$ & 0{,}93 & 0{,}61$^\dagger$ \\
MobileSAM, L3 & 0{,}93 & 0{,}75 & 0{,}97 & 0{,}90 & 0{,}93 & 0{,}75 & 0{,}94 & 0{,}61$^\dagger$ & 0{,}95 & 0{,}86$^\dagger$ & 0{,}95 & 0{,}61$^\dagger$ \\
MobileSAM, L4 & 0{,}78 & 0{,}75 & 0{,}98 & 1{,}61 & 0{,}93 & 1{,}03 & 0{,}82 & 0{,}61$^\dagger$ & 0{,}94 & 1{,}92$^\dagger$ & 0{,}95 & 1{,}63$^\dagger$ \\
ВО, каждый 2-й кадр & 0{,}67 & 0{,}75 & 0{,}92 & 2{,}46$^\dagger$ & 0{,}93 & 1{,}95$^\dagger$ & 0{,}74 & 0{,}85$^\dagger$ & 0{,}97 & 2{,}98$^\dagger$ & 0{,}96 & 2{,}31$^\dagger$ \\
ВО, каждый 4-й кадр & 0{,}51 & 0{,}75 & 0{,}97 & 3{,}15$^\dagger$ & 0{,}96 & 2{,}41$^\dagger$ & 0{,}75 & 0{,}85$^\dagger$ & 0{,}99 & 3{,}46$^\dagger$ & 0{,}97 & 2{,}42$^\dagger$ \\
HM3D (OOD), L0 & --- & --- & --- & --- & --- & --- & 0{,}50 & 0{,}85 & 0{,}50 & 0{,}85 & 0{,}50 & 0{,}85 \\
HM3D (OOD), L3 & --- & --- & --- & --- & --- & --- & 0{,}50 & 0{,}85 & 0{,}50 & 1{,}07 & 0{,}50 & 0{,}82 \\
\bottomrule
\end{tabular}

\end{table}

Одометрия кадр к кадру без замыканий циклов ошибается неравномерно: ATE по сценам от 3,5~см до 3,9~м, медиана 78~см. Ошибка набирается в основном на отдельных шагах --- поворотах на месте у стен. Запас «только метки» молча теряет покрытие: 0,67 для велосипеда и 0,74 для растения (табл.~\ref{tab:variants}). Совместный запас остаётся валидным (0,93 и 0,96), но отказывается в 60 и 86\,\% разбиений. В миссиях планировщик без калибровки находит путь в 23\,\% задач, и 44\,\% этих путей опасны; совместный вариант отказывается в 89\,\% задач. При такой ошибке позы полезной гарантии нет, и метод сообщает об этом вместо ложного сертификата.

Запасы, откалиброванные на ReplicaCAD, мы проверили на сценах HM3D~\cite{hm3d} из OSMa-Bench --- реальных сканах домов в условии \texttt{no\_lights}. Растение видно с траектории лишь в двух из шести одноэтажных сцен. В одной детектор находит его (промах 0,18~м), в другой делит декоративное растение между классами «подставка», «скульптура» и «горшок» (промах 2,9~м). Покрытие 0,5 на двух сценах --- иллюстрация, а не оценка: гарантия не переносится на другое семейство сцен, калибровать нужно там, где работает робот.

\subsection{Качество карты и безопасность}\label{sec:h4}

Для каждой сцены мы сопоставили mIoU карты при точных позах с долей опасных путей планировщика без калибровки. Ранговая корреляция $\rho = -0{,}16$ ($p = 0{,}57$, 15 сцен); для IoU одних запретных классов знак даже обратный ($\rho = +0{,}51$, $p = 0{,}05$). mIoU усредняет по клеткам, а безопасность определяет наихудший объект в сцене (H4).
